\section{官方幻灯片证据链：AI 系统为何必须重写抽象}

本节先沿官方 deck 的顺序还原讲者论证。Yangqing Jia 并不是从“Agent 应该使用哪个框架”出发，而是先去神秘化智能，再分别考察模型、应用、基础设施和硬件/软件共设计。读图时要区分三种内容：可观察的市场或系统证据、讲者基于经历给出的工程判断，以及讲义对 Agent 系统的进一步推论。三者相关，但证据强度不同。

\subsection{去神秘化：复杂结果可以来自可分解机制}

前七页先交代讲者从 researcher、engineer 到 entrepreneur 的路径，再用 Chinese Typewriter Problem 说明一种分析方法：面对看似不可能的复杂输入，不要先诉诸“神秘智能”，而要寻找索引、检索、组合和人机协作机制。slide 6 的最终 build 在巨大字符表上标出常用语义簇，直观展示一个操作者如何把复杂字符空间压缩为可导航结构。这个类比不证明 LLM 与中文打字机机制相同，但它训练了第一性原理的提问方式。

\lecturefigure{slides-images/slide-001.png}{官方 slide 1：系统设计演进与 AI engineer 视角}
\lecturefigure{slides-images/slide-002.png}{官方 slide 2：researcher、engineer、entrepreneur 的经历与开源工作}
\lecturefigure{slides-images/slide-004.png}{官方 slide 4：Chinese Typewriter Problem}
\lecturefigure{slides-images/slide-006.png}{官方 slide 6：字符表上的检索与组合机制最终 build}
\lecturefigure{slides-images/slide-007.png}{官方 slide 7：模型、应用、AI infra 与创业经验的课程地图}

\begin{knowledgebox}{读图：类比的作用是拆问题，不是给出等价证明}
打字机案例中，目标空间极大，但实际语言有频率结构和语义邻近性，工程系统可以利用这些结构降低操作成本。LLM 同样把离散符号映射到连续表示并利用统计规律，但训练方式、泛化机制和自动化程度完全不同。正确结论是“复杂行为值得被分解为数据结构、算法和界面”，而不是“现代模型只是输入法”。类比负责消除神秘感，后续证据仍需来自真实模型和系统。
\end{knowledgebox}

\textbf{课堂提示：}讲者的身份变化很重要。研究者倾向问算法是否新颖，工程师关心系统是否稳定，创业者还要问用户是否付费、供应链是否可得。整讲不断在三种尺度间切换，因此某些强判断是生产经验而非普适定理。阅读时应保留其适用条件，而不是把个人观点抹平。

\subsection{模型浪潮：算法进步、消费增长与 hype 不是同一变量}

slides 9--12 列出 2025 年前后的模型名称、OpenRouter 消费增长、算法阶段类比和 Gartner hype cycle。课程的核心判断是，新算法仍在推动能力边界，推理调用也在增长，所以不能只用“训练融资过热”推断“真实使用没有增长”。另一方面，模型列表与消费曲线也不能证明每家公司、每种商业模式都健康。bubble、adoption 和 unit economics 属于不同层级的问题。

\lecturefigure{slides-images/slide-009.png}{官方 slide 9：GPT、Grok、Gemini、Kimi、DeepSeek、Qwen 与 GPT-OSS}
\lecturefigure{slides-images/slide-010.png}{官方 slide 10：OpenRouter 上持续增长的模型消费}
\lecturefigure{slides-images/slide-011.png}{官方 slide 11：结构创新、MoE、test-time scaling 与 RL 的历史类比}
\lecturefigure{slides-images/slide-012.png}{官方 slide 12：Generative AI 2025 hype cycle}

\begin{knowledgebox}{读图：四张图回答四个不同问题}
slide 9 说明供给侧模型快速迭代；slide 10 说明特定聚合平台上的调用量增长；slide 11 是讲者对算法阶段的历史类比；slide 12 描述市场预期周期。它们共同支持“AI 仍在快速演进”，但不能互相替代。OpenRouter 不是全部市场，hype cycle 不是性能曲线，模型发布日期也不等于商业价值。分析 Agent 产业时必须把 capability、usage、revenue 和 expectation 分开建模。
\end{knowledgebox}

\begin{warningbox}{时间边界}
这些产品、排名和曲线是 Fall 2025 课堂快照。到 2026-08-18，模型版本和市场格局可能已经变化，因此讲义保留它们作为课程论证，不把它们表述为当前排行榜。\textbf{老师强调：}讲者说“看起来没有 bubble”，语境主要指消费需求继续增长；这不是对估值、资本效率或所有 Agent 创业公司的统一结论。
\end{warningbox}

slide 11 的深层价值是把算法创新看成不断改变系统负载的来源。MoE 引入路由与 expert communication，test-time scaling 增加推理阶段的可变计算，RL 需要 rollout 和环境服务。模型层每出现一种新能力，基础设施层都会获得新的调度、内存和可靠性问题，这也是课程从算法自然过渡到系统设计的原因。

\subsection{应用市场：模型能力与产品体验相关，但不等价}

slides 14--19 与 21 把 ToC、prosumer 和 ToB 分开。Elmo 例子强调，优秀 AI companion 需要角色、记忆、交互节奏、内容策略和分发，不是换成更强基础模型就自动获得。消费产品榜单变化快，说明 model improvement 会降低复制门槛并重排入口；prosumer 因工作价值较容易形成付费；enterprise 市场潜力大，但部署、数据、权限和组织采购使落地更慢。

\lecturefigure{slides-images/slide-014.png}{官方 slide 14：产品体验与模型能力相关但相互独立}
\lecturefigure{slides-images/slide-015.png}{官方 slide 15：Elmo AI companion 的产品与用户反馈}
\lecturefigure{slides-images/slide-016.png}{官方 slide 16：消费 AI 产品榜单快速变化}
\lecturefigure{slides-images/slide-017.png}{官方 slide 17：消费产品子集与模型能力提升带来的流动性}
\lecturefigure{slides-images/slide-018.png}{官方 slide 18：prosumer 的明确付费意愿}
\lecturefigure{slides-images/slide-019.png}{官方 slide 19：enterprise AI 的收入增长与不确定入口}
\lecturefigure{slides-images/slide-021.png}{官方 slide 21：企业 AI 平台栈的最终 build}

\begin{knowledgebox}{读图：从模型层到业务价值要经过四道转换}
第一道是能力可用性：模型能否稳定完成任务；第二道是产品体验：交互、记忆和错误恢复是否顺畅；第三道是流程嵌入：权限、数据和现有系统能否接通；第四道是经济性：节省的时间或新增收入是否覆盖推理、集成和运维成本。消费榜单主要暴露前两道，enterprise stack 还要解决后两道。Agent 的价值不应只用 benchmark 或 demo 惊艳度衡量。
\end{knowledgebox}

\begin{importantbox}{ToC、prosumer、ToB 的验证单位}
ToC 常看 retention、session frequency 和分发成本；prosumer 看个人是否愿意为生产率持续付费；ToB 则看流程 KPI、合规、总拥有成本和采购扩张。三类产品可以共享模型与工具，却不能共享一套上线证明。课程中“consumer fluid、enterprise hopeful”表达的是验证节奏不同，不是简单判断哪一侧一定成功。
\end{importantbox}

\textbf{实践经验：}讲者在约 00:32--00:52 强调完美 app experience 与模型相关但独立。对 Agent 团队，这意味着模型升级要通过版本化任务回归，而产品团队仍需拥有状态管理、权限设计、可观测和人机协作。若所有差异都寄托在基础模型，模型供应商下一次升级也可能同时消灭产品护城河。

\subsection{第三支柱：从 scientific compute 到 AI cloud}

slides 23--25 用 Bitter Lesson 和基础设施时间轴建立“第三支柱”论点。scientific computing 面向大规模数值模拟；web service cloud 擅长托管任意代码并搬运网页、图片和视频；data cloud 为 exabyte 级数据处理形成新平台；AI cloud 同时需要高性能异构计算、分布式执行和 cloud-native 开发体验。关键变化是 compute 与 I/O 的比例、硬件同质性和 workload 结构都不同。

\lecturefigure{slides-images/slide-023.png}{官方 slide 23：Bitter Lesson 对可扩展计算方法的强调}
\lecturefigure{slides-images/slide-024.png}{官方 slide 24：scientific、web、data 与 AI cloud 的基础设施谱系}
\lecturefigure{slides-images/slide-025.png}{官方 slide 25：data compute、web services 与 AI compute 对照}

\begin{knowledgebox}{术语消化：I/O、compute 与 AI cloud}
\textbf{I/O（input/output）}是数据在存储、网络和计算设备之间的输入输出搬运；\textbf{compute}是矩阵乘法、归一化等实际数值运算。传统 data processing 常有 I/O 远大于计算，web service 运行任意代码但单请求计算有限；AI training 的矩阵计算极重，同时仍需高速数据与跨 GPU 通信。\textbf{AI cloud}在本讲中不是一个标准化产品名，而是为训练、推理和模型应用共同设计的硬件、调度与开发平台。
\end{knowledgebox}

\begin{knowledgebox}{读图：第三支柱不是“再建一朵云”}
slide 24 的时间轴强调每一代 workload 都催生新的抽象与公司，但旧能力没有消失。AI 平台仍需要对象存储、网络、身份、日志和数据库，只是 GPU 拓扑、数值计算、模型生命周期与供应链成为一等约束。slide 25 中“Compute >> IO”是相对 conventional workload 的简化，不表示 AI 可以忽略数据管道；训练若喂不满 GPU，同样会被 I/O 卡住。
\end{knowledgebox}

\begin{warningbox}{Bitter Lesson 的系统解释边界}
利用更多计算的一般方法长期占优，并不意味着任何低效系统都值得扩张。若 GPU 因数据、通信或故障空转，账面 compute 并没有转化为学习。课程引用 Bitter Lesson 后立即讨论 infra，正是要把“可扩展算法”与“能把硬件变成有效吞吐的系统”放在一起。
\end{warningbox}

\subsection{传统 cloud 价值重写：bare metal 与通用 K8s 的二分都不够}

slides 26--32 把 conventional cloud 的价值拆成 software variety 与 supply-chain flexibility。传统云通过虚拟化让可交换 CPU 资源运行大量不同应用；AI workload 软件栈更集中，却依赖昂贵 GPU、固定互联和大规模训练，live migration 与硬件替换都更困难。slide 28 的供应商/芯片 Venn 图提出抽象层缺口；slide 29 用强烈标题否定 bare metal 与通用 K8s 两个极端；slides 30--31 给出 developer efficiency、GPU reliability、多云供应链、弹性和统一平台等实践目标。

\lecturefigure{slides-images/slide-026.png}{官方 slide 26：传统 cloud 的 software 与 supply-chain 价值命题}
\lecturefigure{slides-images/slide-027.png}{官方 slide 27：conventional cloud 与 AI cloud 的工作负载/供应链对照}
\lecturefigure{slides-images/slide-028.png}{官方 slide 28：云厂商、芯片与 serverless 抽象之间的缺口}
\lecturefigure{slides-images/slide-029.png}{官方 slide 29：bare metal 与通用 Kubernetes 的双重批评}
\lecturefigure{slides-images/slide-030.png}{官方 slide 30：developer efficiency、infra efficiency 与 GPU 故障}
\lecturefigure{slides-images/slide-031.png}{官方 slide 31：多云、弹性、统一训练/推理平台与团队能力}
\lecturefigure{slides-images/slide-032.png}{官方 slide 32：AI infra 不同，但仍遵循经典系统原则}

\begin{knowledgebox}{读图：slide 27 应逐行比较}
Software variety 行说明传统 cloud 要承载数据库、网络、中间件等多样软件，而 AI cloud 更集中于数值计算框架；workload 行说明传统 VM 可承担多种任务，训练作业形状更固定；supply-chain flexibility 行说明虚拟化可迁移 CPU workload，而大型 GPU training 与特定拓扑绑定；interchangeability 行说明不同 GPU/互联并非透明替换。集中软件栈降低部分平台复杂度，却把硬件拓扑和调度复杂度推高。
\end{knowledgebox}

\begin{importantbox}{Locality、elasticity 与 utilization}
\textbf{locality（局部性）}指计算与所需数据、参数或通信伙伴在拓扑上足够接近；\textbf{elasticity（弹性）}指资源能随需求增减；\textbf{utilization（利用率）}描述昂贵设备有多少时间执行有效工作。AI training 希望高 locality 和稳定 gang scheduling，弹性却受 checkpoint、通信组和硬件数量限制。平台设计必须明确优先级，不能假设三者可以无代价同时最大化。
\end{importantbox}

\begin{warningbox}{“K8s is wrong” 的准确语境}
\textbf{老师强调：}讲者不是说 Kubernetes 永远不能管理 AI workload，而是批评把面向无状态微服务的通用抽象原样套到拓扑敏感、长时、gang-scheduled 的 GPU 作业。bare metal 让开发者直接面对设备和故障，通用 K8s 又可能隐藏关键 topology。合理方向是保留 API、调度与生态优势，同时增加 GPU-aware placement、队列、checkpoint、故障恢复和训练/推理原语。
\end{warningbox}

slide 30 将 developer time 与 GPU time 并列，是本组最重要的成本判断。过度追求裸性能会让工程师花大量时间处理环境与故障；过度抽象则可能让 GPU 空转。AI-native platform 的任务是让两种昂贵资源共同高效，而不是只优化某个 dashboard 上的 utilization。

\subsection{回到未来：硬件/软件共设计与 locality 回归}

最后三张教学图没有大段文字，却承担课程标题“back to the future”的视觉论证。slide 34 把早期大型计算设备与现代 GPU server/accelerator box 并置，显示系统从通用分散服务器重新走向高密度专用节点；slide 35 用长期趋势和设备代际暗示计算能力快速增长伴随新的功耗、互联与供应链约束；slide 36 的机场航拍把复杂拓扑具象化：资源数量多不等于任意两个位置之间都能低成本移动。

\lecturefigure{slides-images/slide-034.png}{官方 slide 34：从早期大型计算设备到现代 GPU/accelerator box}
\lecturefigure{slides-images/slide-035.png}{官方 slide 35：计算系统与加速器代际的长期演进}
\lecturefigure{slides-images/slide-036.png}{官方 slide 36：机场拓扑作为 locality 与调度复杂度的视觉类比}

\begin{knowledgebox}{读图：硬件历史图支持的三条结论}
第一，专用化周期会反复出现，通用计算与领域加速器之间不断摆动；第二，单节点变得高密度后，节点内互联、散热和故障域更重要；第三，大规模系统的性能取决于拓扑，而不只是设备总数。图中没有给出严格 benchmark，不能据此比较具体硬件优劣；它的用途是建立系统设计问题：软件抽象必须暴露哪些物理事实。
\end{knowledgebox}

\begin{warningbox}{中心化与边缘化不是一次性选择}
高密度训练倾向中心化，因为需要昂贵互联和电力；隐私、延迟、带宽和可用性又会推动部分推理靠近用户。Agent 系统还可能把规划放在云端、工具执行放在企业网络、轻量模型放在设备。所谓“回到未来”不是回到单一大型机，而是重新承认 locality、容量规划和硬件/软件共设计这些经典约束。
\end{warningbox}

\textbf{课堂提示：}讲者最后说 AI infra 不同但也相同。不同之处是 GPU、模型和拓扑；相同之处是可靠性、可观测、供应链、恢复和组织协作仍决定生产质量。下一部分的综合章节将把这套证据链转成故障恢复、RAG/Agent、调度和平台建设方法。

